\documentclass[12pt]{article}
\usepackage{amsmath,amssymb,amsthm}
\usepackage[margin=1in]{geometry}
\usepackage{longtable,booktabs}

\title{Problem 308: An Amazing Prime-generating Automaton}
\author{Project Euler Solution}
\date{}

\begin{document}
\maketitle

\section{Problem Statement}
Conway's PRIMEGAME is the FRACTRAN program:
\[
    \frac{17}{91},\; \frac{78}{85},\; \frac{19}{51},\; \frac{23}{38},\; \frac{29}{33},\; \frac{77}{29},\; \frac{95}{23},\; \frac{77}{19},\; \frac{1}{17},\; \frac{11}{13},\; \frac{13}{11},\; \frac{15}{2},\; \frac{1}{7},\; \frac{55}{1}
\]

Starting with $n = 2$, repeatedly multiply by the first fraction that yields an integer result. The powers of~2 that appear generate the primes: $2^2, 2^3, 2^5, 2^7, 2^{11}, \ldots$

Find the number of FRACTRAN iterations to produce $2^{p_{10001}}$.

\section{Register Machine Encoding}

The working number's prime factorization encodes registers:
\[
    n = 2^a \cdot 3^b \cdot 5^c \cdot 7^d \cdot 11^e \cdot 13^f \cdot 17^g \cdot 19^h \cdot 23^i \cdot 29^j
\]

Each fraction corresponds to a conditional register operation. The encoded algorithm performs \textbf{trial division} to test primality of each candidate.

\section{Macro-Step Optimization}

The key insight is that the FRACTRAN program's inner loops can be analytically counted.

\subsection{State Representation}
At each ``subtraction round'' entry, the relevant state is $(a_{\mathrm{acc}}, B, C, D)$ with $e=1$.

\subsection{Three Outcomes per Round}
\begin{enumerate}
    \item \textbf{Full round} ($D+B \leq C$): Cost $= 2B + 2(D+B) + 2$ steps. Continue.
    \item \textbf{Overflow} ($D+B > C$, $C > 0$): Cost $= 2B + 2C + 2a_{\mathrm{acc}} + 4$ steps. Next divisor.
    \item \textbf{End of candidate} ($D+B > C$, $C = 0$):
    \begin{itemize}
        \item If $D+B = 1$: candidate is \textbf{prime}.
        \item Otherwise: candidate is \textbf{composite}.
    \end{itemize}
\end{enumerate}

Each macro-step runs in $O(1)$ time, making the simulation practical.

\section{Result}
\[
    \boxed{1539669807660924}
\]

\section{Complexity}
\begin{itemize}
    \item \textbf{Time:} $O\!\left(\sum_{m=1}^{p_{10001}} m\right)$ macro-steps.
    \item \textbf{Space:} $O(1)$.
\end{itemize}

\section{Answer}

\[
\boxed{1539669807660924}
\]

\end{document}
